\section{Executive Summary}

California State University, Long Beach (CSULB), through the Department of Mechanical and Aerospace Engineering, proposes to develop and validate reinforcement learning-augmented multi-model predictive guidance algorithms for hypersonic glide vehicles operating in dynamic threat environments. The proposed research directly supports the NSWC PHD interest in hypersonic development, guidance, navigation and control systems optimized for hypersonic flight, computational modeling, and validation of advanced hypersonic technologies.

The proposed effort builds on prior CSULB-led research by David Daeyoung Lee on model-based numerical predictor-corrector guidance (NPCG) and predictive lateral evasive maneuver guidance for hypersonic glide vehicles capable of avoiding Predicted Intercept Points (PIPs). Unlike conventional no-fly-zone avoidance problems, PIPs are small, continuously updated, and distributed along the vehicle's flight path. This creates a need for online guidance logic that can rapidly evaluate candidate maneuvers, maximize clearance from imminent intercept opportunities, and preserve terminal targeting performance.

Although model-based NPCG can provide accurate entry targeting, it can be sensitive to atmospheric, aerodynamic, and vehicle-model uncertainty. It may also require substantial onboard computation because the guidance computer must repeatedly predict and correct future trajectories during flight. This research will address these limitations by developing multiple NPCG models under different uncertainty assumptions and using their results to train a reinforcement learning agent. The trained RL agent will use sensor-informed state data to perform adaptive guidance actions, seeking to preserve the targeting discipline of physics-based NPCG while improving robustness and real-time computational feasibility.

The expected outcome is a validated adaptive guidance research framework, performance characterization, and transition roadmap to support future NSWC PHD hypersonic GNC research and validation activities.
